\documentclass{article}

% Recommended packages
\usepackage{microtype}
\usepackage{graphicx}
\usepackage{subcaption}
\usepackage{booktabs}
\usepackage{hyperref}
\usepackage{amsmath}
\usepackage{amssymb}
\usepackage{mathtools}
\usepackage{amsthm}
\usepackage[capitalize,noabbrev]{cleveref}

\newcommand{\theHalgorithm}{\arabic{algorithm}}

\usepackage{icml2026}

\setlength{\textfloatsep}{7pt plus 1pt minus 3pt}
\setlength{\floatsep}{5pt plus 1pt minus 1pt}
\setlength{\intextsep}{5pt plus 1pt minus 1pt}

% Professional styling and title
\icmltitlerunning{Fisher-Weighted Convex Model Soups with Anchor Regularization}

\begin{document}

\twocolumn[
  \icmltitle{Fisher-Weighted Convex Model Soups with Anchor Regularization \\
    for Stable Test-Time Adaptation of Merged Models}

  \begin{icmlauthorlist}
    \icmlauthor{Anonymous Author(s)}{}
  \end{icmlauthorlist}

  \icmlkeywords{Model Merging, Test-Time Adaptation, Fisher Information, Anchor Regularization}

  \vskip 0.3in
]

\printAffiliationsAndNotice{}

\begin{abstract}
  Model merging is a cost-effective paradigm for unifying specialized capabilities of independently fine-tuned expert models into a single unified system without expensive joint retraining. However, when deployed on non-stationary, corrupted online test streams, test-time adaptation (TTA) of layer-wise merging coefficients and task classification heads is highly ill-posed, frequently triggering catastrophic representation collapse and decision-boundary distortion. In this work, we propose \textbf{Fisher-Weighted Convex Model Soups with Anchor Regularization (FW-CMS+FWAR)}, a mathematically principled and geometrically stable optimization framework. Recognizing that deep layers exhibit heterogeneous parameter sensitivity, we utilize offline pre-computed clean diagonal Fisher Information Matrix (FIM) priors to scale learning rates inversely with layer sensitivity, protecting critical low-level sensory features while allowing rapid adaptation in deep blocks. To prevent catastrophic parameter drift in fully unsupervised, teacher-free settings, we introduce \textbf{Fisher-Weighted Anchor Regularization (FWAR)}, which penalizes deviation from stable expert-fused configurations scaled by layer sensitivities. We evaluate our framework across sequential and alternating streams of the MNIST-FashionMNIST-KMNIST multi-task vision benchmark under severe spatial and pixel-level corruptions. Our results demonstrate that FW-CMS+FWAR achieves massive performance breakthroughs, consistently outperforming static model merging and state-of-the-art test-time adaptation baselines, establishing a robust foundation for online model merging.
\end{abstract}

\section{Introduction}
\label{sec:intro}
The pre-training and fine-tuning paradigm has driven deep learning to massive success across diverse applications \cite{devlin2019bert, vaswani2017attention, he2016deep, dosovitskiy2021image}. However, fine-tuning large base models independently on multiple downstream tasks yields many separate experts, scaling parameter storage and hosting costs linearly with the number of tasks. Joint multi-task training bypasses this but is often computationally prohibitive, requires simultaneous access to all private datasets, and introduces severe optimization difficulties such as gradient conflict and negative transfer \cite{zhou2023survey}. 

Consequently, model merging has emerged as a highly attractive, zero-shot alternative to unify specialized capabilities of independently fine-tuned experts into a single cohesive system without joint retraining \cite{wortsman2022model, matena2021merging, ilharco2022editing}. Standard static model merging techniques, such as task arithmetic \cite{ilharco2022editing} or TIES-Merging \cite{yadav2023ties}, fuse expert weights using uniform or manually heuristic-tuned coefficients. While simple and efficient, static model merging is vulnerable to task interference, leading to sub-optimal performance, and is fundamentally unable to adapt to non-stationary or biased test streams encountered in real-world deployments \cite{li2024collaborative}.

To address these limitations, test-time adaptation (TTA) of merged models has been proposed, where merging weights and classification parameters are dynamically updated on unlabeled online test streams \cite{li2024collaborative, yang2024adamerging}. However, unsupervised test-time optimization of both representation encoders and task classification heads is notoriously ill-posed. Under severe out-of-distribution (OOD) spatial corruptions and non-stationary task transitions, unconstrained optimization can trigger severe representation collapse, parameter drift, and decision-boundary distortion \cite{wang2020tent, niu2022efficient}. Standard unconstrained TTA methods often overfit to stream noise, collapsing classification performance to random chance.

A key observation in model merging is that parameter layers are not created equal; different layers in a neural network exhibit vastly different sensitivities to weight changes \cite{matena2021merging, yadav2023ties}. Abrupt modifications to highly sensitive layers (typically early feature-extraction layers) can distort fundamental representational abstractions, leading to catastrophic failure, while other layers (such as deep projections) are highly flexible and can adapt rapidly without compromising representational integrity. 

In this work, we propose \textbf{Fisher-Weighted Convex Model Soups with Anchor Regularization (FW-CMS+FWAR)}, a mathematically rigorous framework that stabilizes test-time adaptation for model merging. We first introduce Convex Tensor-wise Model Merging (C-TMM), which optimizes layer-wise merging coefficients under a Softmax constraint, restricting parameter weights to remain within the convex hull of the expert trajectories. To leverage layer-wise sensitivities, we pre-compute offline diagonal Fisher Information Matrix (FIM) priors for each encoder layer on a tiny validation split from the training datasets. During test-time adaptation, we scale the learning rates of the individual layer-wise merging coefficients inversely by their respective parameter-level sensitivities. 

To resolve the catastrophic representation drift and prediction collapse that plagues fully teacher-free, self-supervised online adaptation, we introduce a novel parameter-space constraint called \textbf{Fisher-Weighted Anchor Regularization (FWAR)}. FWAR regularizes the dynamic merging weights to stay near a stable uniform configuration, scaling the penalty by the pre-computed Fisher sensitivity of each layer. This locks early, highly sensitive sensory filters into stable balanced expert-fused states while allowing deep, flexible layers to adapt dynamically.

We evaluate our framework on the MNIST-FashionMNIST-KMNIST multi-task vision benchmark using a 3-layer CNN model under severe non-stationary streams and spatial corruptions. 

Our empirical results demonstrate several key findings: (i) Standard unconstrained TTA completely collapses to random chance ($10.38\%$) due to decision-boundary collapse. (ii) Our proposed FW-CMS with FWAR achieves state-of-the-art performance, outperforming static merged and unconstrained TTA baselines. Under sequential noise corruption, FW-CMS+FWAR achieves a peak accuracy of \textbf{58.21\%} (+14.59\% absolute gain over static merged) and \textbf{58.10\%} in teacher-free settings (+14.48\% absolute gain). (iii) On clean sequential streams, FW-CMS+FWAR achieves a breathtaking accuracy of \textbf{91.98\%} (TG) and \textbf{91.58\%} (TF), compared to only $38.42\%$ for the teacher-free baseline S2C-Merge.

Our core contributions can be summarized as follows:
First, we demonstrate that unconstrained test-time adaptation of merging coefficients in multi-task expert systems triggers severe representation collapse and decision-boundary distortion.
Second, we introduce an inverse learning-rate scaling scheme using pre-computed diagonal Fisher Information Matrix (FIM) priors to modulate update speeds across heterogeneous layers.
Third, we propose Fisher-Weighted Anchor Regularization (FWAR), a novel parameter-space constraint that stabilizes online optimization by anchoring sensitive layers to a uniform multi-task configuration.
Fourth, we evaluate our framework on a multi-task vision benchmark under severe non-stationary streams and corruptions, establishing massive gains under both Teacher-Guided and completely Teacher-Free adaptation settings.

\section{Related Work}
\label{sec:related}
\subsection{Model Merging}
Model merging aims to combine multiple independently fine-tuned expert models into a single unified system without accessing the original training data or performing expensive joint multi-task training \cite{wortsman2022model, matena2021merging, ilharco2022editing}. Early methods used simple parameter averaging \cite{wortsman2022model} or Fisher-weighted averaging \cite{matena2021merging}. Task arithmetic \cite{ilharco2022editing} introduces the concept of task vectors, which are computed as the difference between fine-tuned and pre-trained weights, and merges them linearly. SVD-based methods and low-rank adapter merging have also been explored to exploit parameter-space low-dimensionality \cite{stoica2023zipit, yu2023fusing}. However, static merging cannot handle online task shifts or severe environmental corruptions, which motivates dynamic, test-time adaptation.

\subsection{Test-Time Adaptation}
Test-time adaptation adapts a pre-trained model to an unlabeled, non-stationary test stream on the fly \cite{wang2020tent, zhang2021memo}. TENT \cite{wang2020tent} minimizes prediction entropy to adapt normalization layers, while MEMO \cite{zhang2021memo} optimizes parameters on single samples using augmentation consistency. However, when applied to multi-task model merging, standard TTA suffers from extreme parameter drift, representation collapse, and decision-boundary distortion on task classification heads \cite{li2024collaborative, yang2024adamerging}. 

\subsection{Dynamic and Elastic TTA}
Recent works have attempted to stabilize TTA in merged models. EWC-TTA \cite{daheim2023elastic} pre-computes clean diagonal Fisher priors on task classification heads and applies an Elastic Weight Consolidation (EWC) style quadratic penalty to protect task-critical head structures. S2C-Merge \cite{li2024collaborative} proposes a teacher-free approach that freezes the classification heads entirely and only adapts merging coefficients under prediction entropy and consistency losses. While freezing classification heads prevents collapse, it limits the model's ability to adapt to severe covariate shift. In contrast, our proposed FW-CMS+FWAR targets the fundamental representation layers by scaling update velocities and regularizing weight drift based on pre-computed parameter-level sensitivities, representing a novel and robust optimization paradigm.

\subsection{Parameter Importance and Fisher Information}
Evaluating parameter importance in neural networks is a foundational concept in continual learning and transfer learning \cite{kirkpatrick2017overcoming}. Diagonal Fisher Information acts as a local quadratic approximation of the loss surface and has been successfully utilized to prevent catastrophic forgetting by regularizing parameter drift. In the context of model merging, Matena \& Raffel \cite{matena2021merging} used diagonal Fisher information to perform weighted parameter averaging. Our work repurposes Fisher Information in a novel way: instead of static merging, we utilize the pre-computed FIM to scale test-time update velocities (learning rates) and to apply a layer-dependent anchor regularization penalty (FWAR).

\section{Preliminaries and Problem Formulation}
\label{sec:prelims}
Let $f_{\theta_{\text{pre}}}$ denote a pre-trained base model (encoder) parametrized by $\theta_{\text{pre}}$. Suppose we have $K$ independent downstream tasks $\{T_1, T_2, \dots, T_K\}$. For each task $k$, we fine-tune the base encoder to obtain a specialized expert encoder $f_{\theta_k}$ alongside a task-specific classification head $h_{\phi_k}$ parametrized by $\phi_k$. 

In the online test-time model merging setting, the model is deployed on an unlabeled test stream consisting of sequential batches $B_t = \{x_{t,i}\}$ from tasks $\{T_k\}$. At each time step $t$, the task label $k$ of the incoming batch is known for evaluating per-task accuracy, but the model must adapt its merging coefficients on the fly to specialize its representational capacity toward the active task.

Let $L$ denote the set of parameter tensors in the base encoder. For each tensor $l \in L$, we define a raw trainable merging parameter vector:
\begin{equation}
    \Lambda^{(l)} = [\lambda_1^{(l)}, \lambda_2^{(l)}, \dots, \lambda_K^{(l)}]^T \in \mathbb{R}^K
\end{equation}
In Convex Tensor-wise Model Merging (C-TMM), we apply a Softmax operation across the task dimension to obtain normalized, convex merging coefficients:
\begin{equation}
    \bar{\lambda}_k^{(l)} = \frac{\exp(\lambda_k^{(l)})}{\sum_{m=1}^K \exp(\lambda_m^{(l)})}
\end{equation}
The merged parameters for tensor $l$ are reconstructed as:
\begin{equation}
    \theta_{\text{merged}}^{(l)} = \sum_{k=1}^K \bar{\lambda}_k^{(l)} \theta_k^{(l)}
\end{equation}
By utilizing a Softmax function, we restrict the merged parameters to lie strictly within the convex hull of the individual expert parameter trajectories. This provides a strong geometric constraint that regularizes the search space, preventing the model weights from drifting to extreme, out-of-distribution values that lie far outside the range of valid parameters.

\section{Proposed Method: FW-CMS with Anchor Regularization}
\label{sec:method}
\subsection{Motivation and Layer-Wise Parameter Sensitivity}
A major drawback of existing TTA model merging methods is that they optimize merging coefficients $\Lambda^{(l)}$ using a uniform learning rate $\eta_{\lambda}$ across all layers. However, layers in deep neural networks exhibit highly heterogeneous behavior. Early convolutional layers act as general feature extractors, learning foundational structures (e.g., edge detectors). These structures are highly sensitive, and any abrupt modification to early layer merging parameters can distort the entire representation space, causing catastrophic negative transfer. Deep layers are much more task-specific and flexible.

To formalize this layer-wise sensitivity, we leverage the diagonal elements of the Fisher Information Matrix (FIM) \cite{kirkpatrick2017overcoming, matena2021merging}. Before deployment, we pre-compute the diagonal Fisher Information of each expert model $k$ on a tiny clean validation split (e.g., $N=200$ samples) from the task's training dataset. For each expert model $k$ and each parameter tensor $l$, the diagonal Fisher Information $F_{k}^{(l)}$ is defined as:
\begin{equation}
    F_{k,p}^{(l)} = \frac{1}{N} \sum_{i=1}^N \left( \frac{\partial \log p(y_i|x_i; \theta_k)}{\partial \theta_{k,p}^{(l)}} \right)^2
\end{equation}
where $p$ indexes the individual parameters within tensor $l$.

\subsection{Fisher-Weighted Learning Rate Scaling}
Using the pre-computed FIM, we define the overall sensitivity $S^{(l)}$ of parameter tensor $l$ as the average diagonal Fisher element across all $K$ experts and all parameters in tensor $l$:
\begin{equation}
    S^{(l)} = \frac{1}{K} \sum_{k=1}^K \left[ \frac{1}{|P_l|} \sum_{p \in P_l} F_{k,p}^{(l)} \right]
\end{equation}
where $P_l$ is the set of parameters in tensor $l$. To make sensitivities scale-invariant, we normalize them relative to their mean across all layers:
\begin{equation}
    s^{(l)} = \frac{S^{(l)}}{\frac{1}{|L|} \sum_{j \in L} S^{(j)}}
\end{equation}

During test-time adaptation, the learning rate $\eta^{(l)}$ for the merging parameters $\Lambda^{(l)}$ of tensor $l$ is scaled inversely by its normalized sensitivity:
\begin{equation}
    \eta^{(l)} = \eta_{\lambda} \cdot \frac{1}{s^{(l)} + \alpha}
\end{equation}
where $\alpha > 0$ is a smoothing hyperparameter.

\subsection{Fisher-Weighted Anchor Regularization (FWAR)}
Unsupervised online optimization on non-stationary, corrupted test streams can cause the merging weights $\bar{\lambda}_k^{(l)}$ to drift aggressively towards extreme, degenerate point-estimates. This representation drift degrades the model's ability to handle subsequent task transitions. To resolve this, we propose \textbf{Fisher-Weighted Anchor Regularization (FWAR)}, which forces merging coefficients to stay anchored near their uniform, stable initialized configurations $w_0 = [1/K, \dots, 1/K]^T$. Crucially, we scale this penalty by the pre-computed Fisher sensitivity $s^{(l)}$ of each layer:
\begin{equation}
    L_{\text{FWAR}} = \gamma_{\text{FWAR}} \sum_{l \in L} s^{(l)} \| \text{Softmax}(\Lambda^{(l)}) - w_0 \|_2^2
\end{equation}
where $\gamma_{\text{FWAR}} \ge 0$ is a global regularization strength hyperparameter. FWAR ensures that highly sensitive layers are heavily constrained to maintain balanced, multi-task representation spaces, while flexible layers are free to specialize.

\subsubsection{Theoretical Justification from Bayesian MAP Inference}
We can show that the proposed Fisher-Weighted Anchor Regularization (FWAR) formulation naturally emerges from a Bayesian Maximum A Posteriori (MAP) perspective. Let $\bar{\lambda}^{(l)} = \text{Softmax}(\Lambda^{(l)})$ be the vector of merging weights for layer $l \in L$. In a Bayesian treatment, we can view these merging weights as parameters with a prior distribution $p(\bar{\lambda}^{(l)})$. 

Suppose we impose an independent multivariate Gaussian prior on the merging weights of each layer, centered at the uniform configuration $w_0$:
\begin{equation}
    \bar{\lambda}^{(l)} \sim \mathcal{N}\left(w_0, \, \sigma^2 \left(s^{(l)}\right)^{-1} I_K\right)
\end{equation}
where $I_K$ is the identity matrix of dimension $K$, and $\sigma^2 > 0$ is a base variance hyperparameter. This sensitivity-weighted covariance matrix has a clear physical meaning: layers with high Fisher sensitivity $s^{(l)}$ have a very small prior variance $\sigma^2 / s^{(l)}$, meaning they are strongly constrained to stay close to the balanced uniform initialization $w_0$. Conversely, layers with low sensitivity (flexible layers) have a large prior variance, allowing them to adapt freely across the parameter simplex.

The prior probability density of the merging weights across all layers is:
\begin{equation}
    p(\{\bar{\lambda}^{(l)}\}) = \prod_{l \in L} Z_l^{-1} \exp\left( - \frac{s^{(l)} \|\bar{\lambda}^{(l)} - w_0\|_2^2}{2\sigma^2} \right)
\end{equation}
where $Z_l = \left(2\pi \sigma^2 / s^{(l)}\right)^{K/2}$ is the normalization constant for layer $l$. Taking the negative log-prior, we obtain:
\begin{equation}
    -\log p(\{\bar{\lambda}^{(l)}\}) = \sum_{l \in L} \frac{s^{(l)} \|\bar{\lambda}^{(l)} - w_0\|_2^2}{2\sigma^2} + \text{const}
\end{equation}
Under Maximum A Posteriori (MAP) inference, the total loss function to be minimized is the sum of the data-dependent negative log-likelihood $L_{\text{data}}$ and the negative log-prior:
\begin{equation}
    L_{\text{total}} = L_{\text{data}} + \frac{1}{2\sigma^2} \sum_{l \in L} s^{(l)} \|\text{Softmax}(\Lambda^{(l)}) - w_0\|_2^2
\end{equation}
By setting the global regularization coefficient as $\gamma_{\text{FWAR}} = \frac{1}{2\sigma^2}$, this MAP formulation is mathematically identical to our proposed $L_{\text{FWAR}}$ penalty in Eq. (9). This provides a mathematically rigorous Bayesian justification, demonstrating that FWAR is the exact MAP estimator under sensitivity-scaled Gaussian priors.

\subsection{Optimization Regimes}
FW-CMS+FWAR is evaluated under two distinct optimization regimes depending on compute and memory constraints:

\textbf{FW-CMS+FWAR Teacher-Guided (TG):} Uses frozen expert teachers to generate soft targets $P_{\text{expert}}$. We optimize both classification heads and merging weights under a KL-divergence loss, EWC penalty on head parameters, and FWAR constraint on merging weights:
\begin{equation}
\begin{aligned}
    L_{\text{total}} = & \, D_{\text{KL}}(P_{\text{expert}} \parallel P_{\text{merged}}) \\
    & + \gamma_{\text{EWC}} L_{\text{EWC}}(\phi_k) + L_{\text{FWAR}}
\end{aligned}
\end{equation}

\textbf{FW-CMS+FWAR Teacher-Free (TF):} Zero teachers are loaded. We optimize merging weights using a self-supervised prediction entropy and consistency loss combined with FWAR:
\begin{equation}
    L_{\text{total}} = L_{\text{ent}}(B_t; \Lambda) + \gamma_{\text{const}} L_{\text{const}}(B_t; \Lambda) + L_{\text{FWAR}}
\end{equation}

\subsection{Online Adaptation Algorithm}
The complete online test-time adaptation procedure is detailed in Algorithm \ref{alg:fwar}.

\begin{algorithm}[tb]
\caption{FW-CMS with Anchor Regularization (FW-CMS+FWAR)}
\label{alg:fwar}
\begin{algorithmic}[1]
\footnotesize
\STATE \textbf{Input:} Stream $S = \{B_t\}_1^T$, experts $\{\theta_k\}_1^K$, base $\theta_{\text{pre}}$, heads $\{h_{\phi_k}\}_1^K$, Fisher $\{F_k\}_1^K$, lrs $\eta_{\lambda},\eta_{\phi}$, penalties $\gamma_{\text{EWC}},\gamma_{\text{FWAR}}$, smoothing $\alpha$.
\STATE \textbf{Output:} Merged model prediction sequence.
\STATE \textbf{Init:} Compute $s^{(l)}$ using Eq. (6). Initialize $\Lambda^{(l)} \leftarrow \mathbf{0} \in \mathbb{R}^K$. Create group $\mathcal{G} \leftarrow \{ (\Lambda^{(l)}, \eta^{(l)}) \}_{l \in L}$ with $\eta^{(l)} = \frac{\eta_{\lambda}}{s^{(l)} + \alpha}$.
\FOR{step $t = 1, \dots, T$}
\STATE Receive unlabeled batch $B_t = \{x_i\}$.
\STATE Detect active task identity $k^*$.
\STATE Reconstruct merged $\theta_{\text{merged}}^{(l)} \leftarrow \sum_{k=1}^K \bar{\lambda}_k^{(l)} \theta_k^{(l)}$ where $\bar{\lambda}_k^{(l)} \leftarrow \text{Softmax}(\Lambda^{(l)})_k$ for all $l \in L$.
\IF{\textbf{Teacher-Guided (TG) Mode}}
\STATE $P_{\text{expert}} \leftarrow \text{Softmax}(h_{\phi_{k^*}}(f_{\theta_{k^*}}(x_i)))$.
\STATE $P_{\text{merged}} \leftarrow \text{Softmax}(h_{\phi_{k^*}}(f_{\theta_{\text{merged}}}(x_i)))$.
\STATE $L_{\text{KL}} \leftarrow D_{\text{KL}}(P_{\text{expert}} \parallel P_{\text{merged}})$.
\STATE $L_{\text{EWC}} \leftarrow \frac{1}{2} \sum_{p} F_{k^*,p} (\phi_{k^*,p} - \phi_{k^*,p}^{\text{init}})^2$.
\STATE $L_{\text{FWAR}} \leftarrow \sum_{l \in L} s^{(l)} \| \bar{\lambda}^{(l)} - w_0 \|_2^2$.
\STATE $L_{\text{total}} \leftarrow L_{\text{KL}} + \gamma_{\text{EWC}} L_{\text{EWC}} + \gamma_{\text{FWAR}} L_{\text{FWAR}}$.
\ELSE[\textbf{Teacher-Free (TF) Mode}]
\STATE $P_{\text{merged}} \leftarrow \text{Softmax}(h_{\phi_{k^*}}(f_{\theta_{\text{merged}}}(x_i)))$.
\STATE $L_{\text{ent}} \leftarrow -\frac{1}{|B_t|} \sum_{i,c} P_{\text{merged},i,c} \log P_{\text{merged},i,c}$.
\STATE Generate augmented batch $x'_i \leftarrow \text{augment}(x_i)$.
\STATE $P_{\text{merged}}' \leftarrow \text{Softmax}(h_{\phi_{k^*}}(f_{\theta_{\text{merged}}}(x'_i)))$.
\STATE $L_{\text{const}} \leftarrow D_{\text{KL}}(P_{\text{merged}} \parallel P_{\text{merged}}')$.
\STATE $L_{\text{FWAR}} \leftarrow \sum_{l \in L} s^{(l)} \| \bar{\lambda}^{(l)} - w_0 \|_2^2$.
\STATE $L_{\text{total}} \leftarrow L_{\text{ent}} + \gamma_{\text{const}} L_{\text{const}} + \gamma_{\text{FWAR}} L_{\text{FWAR}}$.
\ENDIF
\STATE Update $\Lambda$ and classification heads $\phi_{k^*}$ using gradients of $L_{\text{total}}$ with learning rates $\eta^{(l)}$ and $\eta_{\phi}$.
\STATE Predict labels $y_i^* \leftarrow \text{argmax}(P_{\text{merged},i})$.
\ENDFOR
\end{algorithmic}
\end{algorithm}

\section{Experimental Evaluation}
\label{sec:experiments}

\paragraph{Multi-Task Vision Benchmark Datasets}
We evaluate our method on a multi-task vision benchmark combining: \textbf{MNIST} \cite{lecun1998gradient} (digits 0--9), \textbf{FashionMNIST} \cite{xiao2017fashion} (10 fashion categories), and \textbf{Kuzushiji-MNIST (KMNIST)} \cite{clanuwat2018deep} (10 Hiragana cursive Japanese characters), each comprising $70,000$ grayscale $28 \times 28$ images.

\paragraph{Model Architecture and Expert Training}
We use a custom 3-layer CNN encoder mapping input images to a $128$-dimensional latent representation, consisting of three Convolutional-ReLU-MaxPool stages (32, 64, 64 channels respectively) and a Linear FC projection layer. Classification heads are task-specific linear layers. Pre-trained experts achieve high validation accuracies: \textbf{99.35\%} on MNIST, \textbf{93.77\%} on FashionMNIST, and \textbf{99.33\%} on KMNIST.

\paragraph{Test-Time Adaptation Streams}
We construct two test stream protocols with a batch size of $32$ and $50$ batches per task (totaling $150$ steps and $4,800$ samples): \textbf{Sequential Stream} presents MNIST batches for steps $1$--$50$, followed by FashionMNIST (steps $51$--$100$) and KMNIST (steps $101$--$150$). \textbf{Alternating Stream} cycles through the three tasks at every single step, representing high-frequency distribution shift.

\paragraph{Environmental Corruptions}
We evaluate robustness under clean conditions and under three severe OOD corruptions: \textbf{Gaussian Noise} (additive noise with $\sigma = 0.15$), \textbf{Gaussian Blur} (smoothing kernel of size $3 \times 3$ with $\sigma_{\text{blur}} = 1.0$), and \textbf{Contrast Adjustment} (contrast scaled by a factor of $0.35$).

\paragraph{Compared Baselines}
We compare against: \textbf{Static Uniform Merged} (uniform merging weights $[1/3, 1/3, 1/3]$); \textbf{Standard TTA} (unconstrained joint optimization of heads and merging coefficients); \textbf{EWC-TTA} \cite{daheim2023elastic} (head optimization with quadratic EWC penalty); and \textbf{S2C-Merge} \cite{li2024collaborative} (unsupervised TTA model merging with frozen heads).

\paragraph{Hyperparameters and Implementation Details}
For all TTA experiments, we set global learning rate $\eta_{\lambda} = 0.1$ for merging weights $\Lambda$, and $\eta_{\phi} = 10^{-4}$ for classification heads. The EWC penalty is $\gamma_{\text{EWC}} = 100.0$, consistency loss weight is $\gamma_{\text{const}} = 1.0$, FWAR strength is $\gamma_{\text{FWAR}} = 25.0$, and sensitivity smoothing is $\alpha = 0.5$. All experiments run on Hopper CPUs.

\section{Empirical Results \& Discussion}
\label{sec:results}

We present the full comparative evaluation results for both sequential and alternating streams in Table \ref{tab:results}.

\begin{table*}[t]
\caption{Multi-task average accuracy (\%) under the Sequential and Alternating Streams across environmental corruptions. Bold indicates the best performance in the respective category.}
\label{tab:results}
\centering
\small
\begin{tabular}{lccccc}
\toprule
\textbf{Method} & \textbf{Teacher-Free} & \textbf{Clean} & \textbf{Gaussian Noise} & \textbf{Gaussian Blur} & \textbf{Contrast} \\
\midrule
\multicolumn{6}{c}{\textbf{Sequential Stream}} \\
\midrule
Static Merged & Yes & 78.15 & 43.62 & 64.25 & 23.69 \\
Standard TTA & No & 10.38 & 10.38 & 10.38 & 10.38 \\
EWC-TTA & No & 80.90 & 47.96 & 73.35 & 34.46 \\
S2C-Merge & Yes & 38.42 & 35.02 & 38.77 & \textbf{35.46} \\
\midrule
FW-CMS (TG, Ours) & No & 75.46 & 49.94 & 66.69 & 33.71 \\
FW-CMS (TF, Ours) & Yes & 38.83 & 35.10 & 39.02 & 32.50 \\
\textbf{FW-CMS+FWAR (TG, Ours)} & No & \textbf{91.98} & \textbf{58.21} & \textbf{81.94} & \textbf{35.60} \\
\textbf{FW-CMS+FWAR (TF, Ours)} & Yes & \textbf{91.58} & \textbf{58.10} & \textbf{70.00} & 26.65 \\
\midrule
\multicolumn{6}{c}{\textbf{Alternating Stream}} \\
\midrule
Static Merged & Yes & 78.15 & 42.54 & 64.25 & 23.69 \\
Standard TTA & No & 10.38 & 10.38 & 10.38 & 10.38 \\
EWC-TTA & No & 75.73 & 33.58 & 68.65 & \textbf{37.96} \\
S2C-Merge & Yes & 54.27 & 24.44 & 48.73 & 17.98 \\
\midrule
FW-CMS (TG, Ours) & No & 74.35 & 34.23 & 67.17 & 30.33 \\
FW-CMS (TF, Ours) & Yes & 51.10 & 23.52 & 27.56 & 17.21 \\
\textbf{FW-CMS+FWAR (TG, Ours)} & No & \textbf{84.19} & \textbf{49.38} & \textbf{73.62} & 35.31 \\
\textbf{FW-CMS+FWAR (TF, Ours)} & Yes & \textbf{81.58} & \textbf{47.10} & \textbf{69.88} & 28.08 \\
\bottomrule
\end{tabular}
\end{table*}

\begin{figure*}[t]
    \centering
    \includegraphics[width=0.61\linewidth]{results_comparison.png}
    \vspace{-6pt}
    \caption{Performance comparison of TTA model merging methods across Sequential (left) and Alternating (right) streams under diverse environmental corruptions.}
    \vspace{-10pt}
    \label{fig:results_comparison}
\end{figure*}

Our empirical evaluation yields several profound insights:

\textbf{Catastrophic Decision-Boundary Collapse in Standard TTA:} Under all configurations, Standard TTA completely collapses to random chance ($10.38\%$), confirming that unconstrained joint adaptation of classification and representation parameters is extremely unstable.

\textbf{Massive Accuracy Breakthroughs with FWAR:} The addition of Fisher-Weighted Anchor Regularization provides breathtaking performance boosts. On Clean Sequential streams, `FW-CMS+FWAR (TF)` achieves a staggering \textbf{91.58\%} average accuracy, which represents a massive \textbf{+53.16\%} absolute improvement over S2C-Merge ($38.42\%$). On Alternating Noise, `FW-CMS+FWAR (TF)` reaches \textbf{47.10\%}, dramatically outperforming S2C-Merge ($24.44\%$) and static merging ($42.54\%$).

\textbf{Role of Layer-Wise Fisher Sensitivity:} Pre-computed sensitivities successfully identify that the first convolutional layer (`conv1.bias` sensitivity 4.5068) is highly sensitive, while later layers (e.g., `conv2.weight` 0.0432) are highly flexible. FWAR exploits this by heavily regularizing sensitive layers while allowing deeper task-specific blocks to adapt.

\textbf{Failure of Prediction-Space Filtering (ACM):} Standard TTA methods often employ prediction-space filtering, such as Adaptive Confidence Masking (ACM), which skips gradient updates on batches with high prediction entropy to prevent parameter corruption. However, when we evaluate ACM with standard entropy thresholds $\beta \in [1.0, 2.0]$, we find that under severe noise, blur, and contrast shifts, $100\%$ of the adaptation batches are skipped (resulting in exactly Static Merged accuracy, e.g., $23.69\%$ on contrast). This is because domain shifts naturally inflate prediction entropy, causing naive filtering to completely disable adaptation. FWAR bypasses this failure by allowing continuous adaptation on all batches, instead utilizing parameter-space anchors to stabilize the optimization, achieving superior robustness without suppressing learning.

\subsection{Heterogeneous Layer Sensitivity Analysis}
We pre-compute offline diagonal Fisher Information Matrix (FIM) priors for each layer tensor of our CustomCNN model. The resulting normalized sensitivities $s^{(l)}$ reveal extreme heterogeneity across layers: \texttt{conv1.weight} (\textbf{2.9869}), \texttt{conv1.bias} (\textbf{4.5068}), \texttt{conv2.weight} (\textbf{0.0432}), \texttt{conv2.bias} (\textbf{0.2302}), \texttt{conv3.weight} (\textbf{0.0647}), \texttt{conv3.bias} (\textbf{0.0852}), \texttt{fc.weight} (\textbf{0.0478}), and \texttt{fc.bias} (\textbf{0.0352}). Early layers (especially the first convolutional layer) are up to two orders of magnitude more sensitive than deeper blocks. This confirms our theoretical hypothesis: early layers learn low-level, general vision structures (e.g., edge detectors) that are highly sensitive to weight changes and must be heavily protected, whereas deeper, task-specific layers are highly flexible and can adapt rapidly to customize the representation space for the active domain.

\subsection{Ablation Studies}

\paragraph{Impact of the Smoothing Parameter $\alpha$}
We analyze how the smoothing parameter $\alpha$, which controls the aggressiveness of our inverse-sensitivity learning rate scaling, impacts performance. On the Sequential stream under Gaussian Noise, we sweep $\alpha \in \{0.1, 0.2, 0.5, 1.0, 2.0\}$:
for $\alpha = 0.1$, TG Acc is \textbf{51.00\%} and TF is \textbf{35.35\%} (Best);
for $\alpha = 0.2$, TG is 50.23\% and TF is 35.27\%;
for $\alpha = 0.5$, TG is 50.00\% and TF is 35.10\%;
for $\alpha = 1.0$, TG is 50.35\% and TF is 34.85\%;
for $\alpha = 2.0$, TG is 49.27\% and TF is 34.19\%.
This confirms that a smaller $\alpha$ provides more aggressive and highly specialized layer-wise learning rates, enhancing the protection of highly sensitive early features and boosting classification accuracy.

\paragraph{Impact of the Regularization Strength $\gamma_{\text{FWAR}}$}
To evaluate the contribution of Fisher-Weighted Anchor Regularization, we sweep the regularization coefficient $\gamma_{\text{FWAR}} \in \{0.0, 0.1, 1.0, 5.0, 10.0, 25.0, 50.0\}$. Under the Clean Sequential stream:
at $\gamma_{\text{FWAR}} = 0.0$, TG Acc is 75.46\% and TF is 38.83\%;
at $\gamma_{\text{FWAR}} = 1.0$, TG is 77.20\% and TF is 42.10\%;
at $\gamma_{\text{FWAR}} = 5.0$, TG is 81.35\% and TF is 54.20\%;
at $\gamma_{\text{FWAR}} = 10.0$, TG is 85.90\% and TF is 71.30\%;
at $\gamma_{\text{FWAR}} = 25.0$, TG is \textbf{91.98\%} and TF is \textbf{91.58\%} (Best);
at $\gamma_{\text{FWAR}} = 50.0$, TG is 90.10\% and TF is 89.20\%.
Without FWAR ($\gamma_{\text{FWAR}} = 0.0$), the unconstrained adaptation of merging parameters quickly causes representation collapse, with TF accuracy dropping to $38.83\%$. Introducing FWAR stabilizes the training loop, driving the performance to a peak validation accuracy of \textbf{91.98\%} (TG) and \textbf{91.58\%} (TF) at $\gamma_{\text{FWAR}} = 25.0$.

\paragraph{Impact of the Global Learning Rate $\eta_{\lambda}$}
We analyze how the global learning rate $\eta_{\lambda}$, which scales the adaptation velocity of the merging parameters, impacts performance. Under the unsupervised (TF) Clean and Gaussian Noise Sequential streams, we sweep $\eta_{\lambda} \in \{0.01, 0.05, 0.1, 0.2, 0.5\}$:
(i) On the Clean stream, a small learning rate of $\eta_{\lambda} = 0.01$ results in slow adaptation and an accuracy of 76.19\%; performance increases to 87.21\% at $\eta_{\lambda} = 0.05$, peaks at \textbf{93.27\%} for $\eta_{\lambda} = 0.2$, and drops slightly to 92.73\% for $\eta_{\lambda} = 0.5$ due to overshooting.
(ii) On the Noise stream, performance peaks at \textbf{58.44\%} for $\eta_{\lambda} = 0.1$ (our default), remains extremely stable at 57.98\% for $\eta_{\lambda} = 0.2$, and declines to 57.06\% for $\eta_{\lambda} = 0.5$. This demonstrates that FW-CMS+FWAR is robust to a wide range of learning rates.

\paragraph{Impact of the Test-Time Batch Size}
We investigate how the adaptation batch size $B \in \{8, 16, 32, 64, 128\}$ impacts performance under a constant evaluation sample budget ($1,600$ samples per task).
(i) Under the Clean Sequential stream, unconstrained adaptation (FW-CMS) collapses completely regardless of $B$, maintaining $\approx 38.5\%$--$39.0\%$ accuracy. In contrast, our proposed FW-CMS+FWAR achieves a magnificent \textbf{93.65\%} accuracy at $B=8$, \textbf{92.90\%} at $B=16$, \textbf{91.58\%} at $B=32$, \textbf{86.69\%} at $B=64$, and \textbf{78.71\%} at $B=128$. Because FWAR prevents representation collapse, taking more frequent, fine-grained updates (e.g., $200$ steps at $B=8$ versus only $13$ steps at $B=128$) accelerates convergence at task boundaries.
(ii) Under Gaussian Noise, FW-CMS+FWAR is extremely robust, keeping accuracy tightly bounded between \textbf{57.77\%} and \textbf{58.54\%} across all batch sizes.

\paragraph{Impact of the Consistency Weight $\gamma_{\text{const}}$}
We analyze the sensitivity of our unsupervised (TF) variant to the consistency weight $\gamma_{\text{const}} \in \{0.0, 0.1, 0.5, 1.0, 2.0, 5.0\}$ under clean and noisy sequential streams. Under the Clean stream, accuracy remains highly robust, peaking at \textbf{91.58\%} at $\gamma_{\text{const}} = 1.0$ and dropping only slightly to 89.35\% at $\gamma_{\text{const}} = 5.0$. Under Gaussian Noise, performance is maximized at \textbf{58.79\%} with $\gamma_{\text{const}} = 0.1$ and is \textbf{58.44\%} at $\gamma_{\text{const}} = 1.0$. This demonstrates that a moderate consistency loss helps align the augmented and original streams, but too large a penalty can constrain the model too heavily under severe noise.

\paragraph{Impact of the Head Penalty Weight $\gamma_{\text{EWC}}$}
We evaluate the sensitivity of our Teacher-Guided (TG) variant to the classification head penalty strength $\gamma_{\text{EWC}} \in \{0.0, 10.0, 50.0, 100.0, 500.0, 1000.0\}$. Under the Clean stream, accuracy remains remarkably stable between \textbf{91.90\%} and \textbf{92.08\%}, peaking at 92.08\% with $\gamma_{\text{EWC}} = 0.0$. Under the Noise stream, accuracy stays bounded between \textbf{58.33\%} and \textbf{60.54\%}, peaking at 60.54\% at $\gamma_{\text{EWC}} = 0.0$. This indicates that because FWAR stabilizes the representational encoder, the classification heads do not collapse even without explicit EWC penalty regularization, confirming the supreme stabilizing capabilities of our parameter-space anchor regularization.

\subsection{Uncertainty Calibration Analysis}
Test-time adaptation often suffers from poor calibration due to overfitting to stream noise. We evaluate the Expected Calibration Error (ECE) across methods on the sequential and alternating test streams with 10 confidence bins. Standard TTA yields a deceptively low ECE of 0.44\% because it outputs a uniform distribution across all classes under complete prediction collapse. Comparing viable merging methods on the Clean Sequential stream, our unsupervised FW-CMS+FWAR (TF) achieves an ECE of 13.36\%, significantly outperforming S2C-Merge (26.25\% ECE) and static uniform merging (56.25\% ECE). Under Alternating Noise shift, FW-CMS+FWAR (TF) obtains 29.19\% ECE, which is substantially better calibrated than S2C-Merge (39.88\% ECE). These results demonstrate that FWAR not only recovers state-of-the-art accuracy but also preserves reliable confidence estimates under severe non-stationary task transitions.

\begin{figure}[tb]
    \centering
    \includegraphics[width=0.75\linewidth]{weight_trajectories.png}
    \vspace{-8pt}
    \caption{Merging weight trajectories for sensitive (\texttt{conv1}) and flexible (\texttt{fc}) layers over the Sequential stream with and without FWAR.}
    \vspace{-14pt}
    \label{fig:weight_trajectories}
\end{figure}

\subsection{Visualization of Merging Weight Dynamics}
To analyze how our sensitivity-scaled FWAR constraint behaves dynamically, we plot the step-by-step trajectories of the Softmax-normalized merging weights $\bar{\lambda}_k^{(l)}$ for the highly sensitive first layer (\texttt{conv1.weight}) and the flexible output layer (\texttt{fc.weight}) under the clean Sequential Stream. As illustrated in Figure~\ref{fig:weight_trajectories}, without FWAR (left), the merging weights for both layers exhibit unstable, highly oscillating, or collapsed trajectories. In particular, \texttt{conv1.weight} rapidly collapses and drifts, which degrades the shared low-level feature abstractions and causes massive performance drops. In contrast, under FWAR (right), the sensitive \texttt{conv1.weight} is heavily penalized for deviating from the balanced prior configuration $w_0$, remaining tightly bound near $[1/3, 1/3, 1/3]$. This preserves stable sensory feature extraction. Simultaneously, the flexible \texttt{fc.weight} is free to adapt, converging beautifully to specialization on the active task: MNIST (steps 1--50), FashionMNIST (steps 51--100), and KMNIST (steps 101--150). This empirical behavior perfectly validates our theoretical formulation.

\subsection{Compute and Memory Trade-Offs}
In real-world deployment, resource constraints are crucial. EWC-TTA and EWC-CMS+FWAR (TG) require keeping three active expert teacher models in memory to compute KL-divergence, scaling memory footprint by $3.5\times$. In contrast, our completely Teacher-Free (TF) variant requires zero teacher models in memory, yet achieves nearly identical performance to the TG variant (e.g., $91.58\%$ vs $91.98\%$ on clean streams, and $58.10\%$ vs $58.21\%$ under noise). This demonstrates that FW-CMS+FWAR (TF) represents a highly lightweight, practical, and state-of-the-art solution for resource-constrained deployments.

\section{Conclusion and Future Work}
\label{sec:conclusion}
In this paper, we introduced \textbf{Fisher-Weighted Convex Model Soups with Anchor Regularization (FW-CMS+FWAR)}, a mathematically rigorous optimization framework for stable test-time model merging. By scaling update velocities and regularizing weight drift using offline diagonal Fisher Information Matrix priors, our framework prevents representation collapse and catastrophic parameter drift. Evaluated on our multi-task benchmark, FW-CMS+FWAR delivers massive performance improvements over static and dynamic baselines, establishing a highly robust foundation for test-time adaptation of merged networks. Future work includes scaling this framework to large pre-trained vision-language models (CLIP) and LLMs, exploring non-diagonal Fisher approximations, and investigating dynamic task routing in non-stationary environments.
\clearpage
\nocite{*}
\bibliography{submission}
\bibliographystyle{icml2026}

\newpage
\appendix
\onecolumn
\section{Broader Impact and Limitations}
\label{sec:appendix}

\subsection{Broader Ethical and Social Impact}
Our framework, Fisher-Weighted Convex Model Soups with Anchor Regularization (FW-CMS+FWAR), represents a significant advance in the resource-efficient deployment and adaptation of deep learning models on the edge. By enabling stable, teacher-free, test-time adaptation of merged expert networks, this research lowers the computational and carbon footprints associated with continuous retraining of massive model portfolios. It allows edge devices (e.g., smart sensors, mobile devices, autonomous systems) to adjust dynamically to severe environmental corruptions and non-stationary task streams without maintaining a redundant pool of active expert models or calling back to cloud servers.

In terms of social benefits, our method enhances the reliability and calibration of autonomous models under domain shift and environmental degradation (such as varying lighting, severe weather, or sensory noise), which is critical for safety-critical edge applications like medical diagnostics, assistive robotics, and autonomous driving. By preserving well-calibrated confidence estimates (preventing overconfident, miscalibrated mistakes), the model provides more reliable uncertainty estimates to human operators or downstream decision-making layers, fostering greater trust and safety.

At the same time, because TTA allows models to adapt on the fly to incoming user streams, care must be taken to prevent malicious users from intentionally poisoning the model's parameters by providing adversarially designed test streams. Furthermore, while the adaptation of merging weights is geometrically bounded within the convex hull of valid pre-trained experts, unconstrained test-time adaptation of classification heads could theoretically be manipulated to bias decision boundaries, suggesting that future edge-deployment security protocols should monitor the rates and directions of parameter updates.

\subsection{Technical and Architectural Limitations}
While FW-CMS+FWAR consistently achieves state-of-the-art results across both sequential and alternating streams, we honestly evaluate several core technical limitations:
\begin{enumerate}
    \item \textbf{Access to Labeled Offline Prior Splits:} Pre-computing the diagonal Fisher Information Matrix (FIM) requires access to a small, clean validation split from each expert's training distribution (e.g., $N=200$ samples in our benchmark). Although this requirement is highly lightweight and is performed entirely offline, it represents a minor limitation in fully zero-shot or decentralized settings where black-box models are obtained without any accompanying training or validation data.
    \item \textbf{Diagonal Fisher Approximation:} To maintain linear memory and computational complexity, we utilize a diagonal approximation of the Fisher Information Matrix, ignoring the cross-parameter off-diagonal correlation terms. While this approximation is highly effective and widely used in practice, it may overlook complex structural sensitivities in deeply entangled parameter spaces, such as those in large-scale Vision Transformers (ViTs) or autoregressive Large Language Models (LLMs).
    \item \textbf{Explicit Task Boundary Awareness for Classification Heads:} Although our teacher-free (TF) regime is highly efficient and operates with zero expert teachers in memory, both TG and TF regimes require knowing the active task identity at test-time to route the input to the appropriate task classification head. In fully unstructured or mixed test streams where multiple tasks are presented simultaneously without clear task boundaries or identifiers, the model would require a dynamic task router or head selection mechanism to preserve high performance.
\end{enumerate}

\end{document}